% backmatter/appendixE.tex —— 附录 E 术语、符号与方法对照
% 纪律：本附录不引入新术语。术语以 spec/glossary.md 为唯一真源，符号以 frontmatter/notation.tex 为准，
% 方法的六问定位逐条来自 research/cards/*.md（每张卡的「六问定位」段）。
% 本附录不参与命题判决，也不产生任何新数字。
\chapter{术语、符号与方法对照}
\label{app:crossref}

% ---------------- E.1 ----------------
\section{为什么需要一张对照表}

\begin{failurebox}
同一个词在不同论文里可能指不同对象，直接比较容易误读。
举一个具体的：Reflexion 说它“有记忆”\cite{reflexion-2303}，
MEM1 说它“有记忆”\cite{mem1-2506}，Mem-$\alpha$ 也说它“有记忆”\cite{memalpha-2509}。
三者的对象并不相同：Reflexion 的记忆是供后续 trial 读取的自然语言反思，实践中保留近期条目；是否跨新任务须核对复位边界；
MEM1 区分 PPO 训练的策略参数与任务内生成的文本内部状态；
Mem-$\alpha$ 用 RL 训练记忆构建策略，运行时该策略操作记忆内容，架构与参数是不同对象。
把三者合并成一句“加了记忆就能跨任务变好”，会掩盖反思文本、任务内状态与记忆构建策略的差别。
比较前须先确定更新对象和复位边界，再讨论第 5 章 §\ref{sec:ch05-forms} 所列的记忆形式。
\end{failurebox}

\begin{sloppypar}
因此本附录做三件事：把本书用到的术语压成一句话定义并点出它最容易混的对象（§\ref{sec:appE-terms}）；
把全书符号集中列出并标明主要讨论章节（§\ref{sec:appE-symbols}）；
把六问记录中的更新对象、反馈与时间尺度压缩成方法索引，使“改什么 / 根据什么改 / 落在哪一层”三栏可以横向比较（§\ref{sec:appE-methods}）。
\end{sloppypar}

\begin{principlebox}
对照表帮助核对比较对象。把两篇论文的数字放在一起前，先确认两处术语所指的对象、时间尺度与结果单位。若不同，先说明比较的是完整方案还是某个机制；共同任务与资源下可以比较整体效果，但不能由此直接归因于某一组件。
\end{principlebox}

% ---------------- E.2 ----------------
\section{术语对照}
\label{sec:appE-terms}

下表是 \texttt{spec/glossary.md} 主表的压缩版，只留三项：本书定义、最容易混的对象、主要讨论章节。
完整定义与不许使用的别名一律以术语表为准。

\begingroup
\begin{longtable}{>{\raggedright\arraybackslash}p{2.4cm}>{\raggedright\arraybackslash}p{6.0cm}>{\raggedright\arraybackslash}p{3.6cm}>{\raggedright\arraybackslash}p{1.4cm}}
\caption{核心术语：定义、常见混淆与主要讨论章节}\label{tab:appE-terms}\\
\toprule
术语 & 本书定义（一句话） & 最容易混的对象 & 章节 \\
\midrule
\endfirsthead
\multicolumn{4}{l}{\tablename\ \thetable（续）}\\
\toprule
术语 & 本书定义（一句话） & 最容易混的对象 & 章节 \\
\midrule
\endhead
\midrule
\multicolumn{4}{r}{续下页}\\
\endfoot
\bottomrule
\endlastfoot
步 / 回合 / 版本 & $t$ 编号回合内的步，$e$ 编号一次任务尝试，$k$ 标识所用系统版本。该次轨迹记为 $\tau_{k,e}$ &
把上下文增长直接当作跨任务学习；须检查保存与读取边界 & 1 \\
\midrule
自进化 & 依据经验与反馈修改持久组件，且该变化可被独立检验 &
自精炼（只改善当前输出，系统无持久变化） & 1 \\
\midrule
持久组件 & $\theta$ 参数、$P$ 指令与持久策略、$M$ 经验记忆、$L$ 工具与技能库、$G$ 工作流与角色结构 &
不限定分量、只说“长期存在的状态”的说法 & 1 \\
\midrule
系统版本 & $A_k=(\theta,P,M,L,G;\mathcal V)$，一个可发布的系统状态 &
把“配置”一词当作同义词；本书的$A_k$包含评价设置，$\mathcal V$默认冻结 & 1 \\
\midrule
验证器 & 任务划分、验收器与最终测试集的合称，\textbf{默认冻结} &
只指“打分模型”；改动任一项都须登记评价版本与可比范围 & 1 \\
\midrule
可更新集合 & 允许被修改的组件集合，每扩大一档都要重答“凭什么认为有用” &
把授权更新范围当成实际权限；接口约束仍需执行隔离 & 2 \\
\midrule
修改器 & $(\mathcal S,H)$：$\mathcal S$ 是改哪个组件，$H$ 是怎么选 &
把“优化器”当作全部；本书要求两者分开写 & 2 \\
\midrule
配对比较 & 同一道题上比较两个版本，以题为单位取差 &
均值差可指向同一目标量，但配对与独立设计的方差口径不同 & 3 \\
\midrule
封存集 & 不参与本轮候选产生与选择的最终测试集，终评前冻结候选 &
终评结果若用于再选版本，就须登记适应性使用 & 3 \\
\midrule
任务族 & 共享一类失败模式与解法的分组；跨新族推断须匹配族级设计 &
把族名当独立性证明，或断言族级区间一定更宽 & 3 \\
\midrule
接地 & 具有独立环境、执行器或外部规则依据的验证信号 &
模型自我赞同；执行检查的有效性仍须核验 & 4、14 \\
\midrule
拟合目标 & 持久物被拟合到求解器（当前模型的失败分布）还是拟合到问题域（与谁读它无关的事实） &
把拟合目标当成迁移保证；两者都须在新读取者上检验 & 5 \\
\midrule
反号失效 & 给定任务与各自对照下，持久物收益在换模型或换 harness 后由正变负 &
E03c计数为逐题点估计，不是单条经验的因果归因 & 5、15 \\
\midrule
记忆操作 & 对 $M$ 的四类动作：写入、检索、淘汰、重写 &
“记忆管理”这种不区分动作的泛指 & 6 \\
\midrule
演化成本 & 产生并验证一次更新消耗的全部资源 &
只算训练算力的账；构建、验证、回归三部分都要进 & 3、8、13 \\
\midrule
回本点 & 演化成本被后续任务的推理节省抵消所需的任务次数 &
“总成本除以单次节省”；不给 $\Delta$cost 就算不出 & 8、13 \\
\end{longtable}
\endgroup

\begin{sloppypar}
这批术语用于提出\textbf{可检查的对象问题}。
给定一个系统，你可以逐条回答“它改了 $A_k$ 的哪个分量”“反馈来自哪一层”“用哪个集检验”。
缺少必要定义或运行信息时，应保留未核查状态，不能根据方法名补出答案。
\end{sloppypar}

% ---------------- E.3 ----------------
\section{符号对照}
\label{sec:appE-symbols}

\begin{sloppypar}
全书符号见表 \ref{tab:appE-sym}。使用规则有三条：
主线符号按前置符号表使用；外部公式与局部例题另有定义时必须明示，不能直接混接；
带帽量（如 $\widehat\Delta$）报告点估计及其不确定性。若设计不足以支持有效区间，应说明原因，不强行补一个看似精确的区间。
\end{sloppypar}

\begin{longtable}{>{\raggedright\arraybackslash}p{3.0cm}>{\raggedright\arraybackslash}p{5.4cm}>{\raggedright\arraybackslash}p{1.2cm}>{\raggedright\arraybackslash}p{3.8cm}}
\caption{符号：含义、主要章节与常见混淆}\label{tab:appE-sym}\\
\toprule
符号 & 含义 & 章节 & 常犯的错误 \\
\midrule
\endfirsthead
\multicolumn{4}{l}{\tablename\ \thetable（续）}\\
\toprule
符号 & 含义 & 章节 & 常犯的错误 \\
\midrule
\endhead
\midrule
\multicolumn{4}{r}{续下页}\\
\endfoot
\bottomrule
\endlastfoot
$t$ / $e$ / $k$ & 步 / 回合 / 版本 & 1 & 用“轮次”指 $t$，用“轨迹”指 $e$ \\
$H_{k,e,t}$ / $\tau_{k,e}$ & 工作上下文 / 单次任务记录 & 1 & 把上下文长度当作记忆容量 \\
$A_k$ / $\widetilde A_{k+1}$ & 系统版本 / 候选版本（未过 Gate） & 1、2 & 把候选当结论报告 \\
$\theta,P,M,L,G$ & 五类持久组件 & 1 & 漏掉 $G$，于是结构变化无处登记 \\
$\mathcal V$ & 验证器（默认冻结） & 1 & 把它写进 Updatable 而不再回答检验问题 \\
Updatable & 可更新集合 & 2 & 扩大集合却沿用旧结论 \\
$\mathcal S,H$ & 修改器的搜索空间与优化算法 & 2 & 只报算法不报搜索空间 \\
$Q$ & 选择验证集的查询预算 & 3 & 把选择验证集当封存集用 \\
$\widehat\Delta_k(b)$ & 预算 $b$ 下的收益估计 & 3 & 只报点估计不报区间 \\
 $FGT_t$ / $BWT_t$ & 遗忘量 / 后向迁移 & 16 & 把两者当成同一个量的正负号 \\
$CPG_t$ / $TP$ & 单位增益成本 / 吞吐 & 3 & 不写单位就做跨论文比较 \\
$\Delta_{\text{migrate}}$ & 换模型或换 harness 后的迁移收益 & 5 & 不跑换读取者的静态基线就下结论 \\
 $C_{\text{evolve}}$ / $N_{\text{break-even}}$ & 演化成本 / 回本点 & 8、13 & 只算训练算力 \\
$k^*$ & 技能固化的临界重复次数 & 8 & 把“越练越好”当默认 \\
\end{longtable}

% ---------------- E.4 ----------------
\section{方法对照}
\label{sec:appE-methods}

\begin{sloppypar}
下表按更新对象组织定位；它是完整六问记录的压缩索引，不是效果或自主性排名。
读者可从 \path{references/sources.csv} 按方法定位文献版次，再查书末对应引文。\path{research/cards/} 收录了部分方法的六问记录，例如 Voyager。完整维护档案另有 \path{agent/source-map.md}，记录本书主张与来源的位置对应。
\end{sloppypar}

\begingroup
\footnotesize
\begin{longtable}{>{\raggedright\arraybackslash}p{2.9cm}>{\raggedright\arraybackslash}p{2.4cm}>{\raggedright\arraybackslash}p{1.9cm}>{\raggedright\arraybackslash}p{1.5cm}>{\raggedright\arraybackslash}p{3.9cm}}
\caption{方法定位：更新对象、反馈与时间尺度}\label{tab:appE-methods}\\
\toprule
方法 & 改什么 & 主反馈 & 层 & 本书中的角色 \\
\midrule
\endfirsthead
\multicolumn{5}{l}{\tablename\ \thetable（续）}\\
\toprule
方法 & 改什么 & 主反馈 & 层 & 本书中的角色 \\
\midrule
\endhead
\midrule
\multicolumn{5}{r}{续下页}\\
\endfoot
\bottomrule
\endlastfoot
ReAct\cite{react-2210} & 此处提示式设置只更新任务内上下文；原文另有微调实验 & 环境观测 & $t$ & 执行基线；方法名不能代替状态边界检查 \\
Self-Refine\cite{self-refine-2303} & 不改 $A_k$，只在 $e$ 层内改输出 & 模型自评（无外部接地） & $e$ & 第 4 章“内在信号”一端的对照 \\
Reflexion\cite{reflexion-2303} & 反思条目；跨 trial 保留 & 外部或内部反馈 + 反思 & 按任务边界判定 & 同题重试保留不自动构成跨新任务更新 \\
A-Mem\cite{amem-2502} & $M$ 的内容与链接集合 & 模型生成的属性与关系判断 & $k$（$M$） & 模型参与组织链接；第6章 \\
VOYAGER\cite{voyager-2305} & $L$（技能库，可执行代码） & 环境、执行错误与模型自评 & $k$（$L$） & 技能复用实例；未测定本书的成本临界值 \\
DSPy\cite{dspy-2310} & 模块指令、示例等参数；程序结构另辨 & 用户指定评价指标 & $k$、离线 & 编译期优化的原型；第 7 章 \\
AFLOW\cite{aflow-2410} & $G$（节点与边）为主 & 任务得分与搜索经验 & $k$、离线 & 工作流搜索；第 9 章 \\
DGM\cite{dgm-2505} & 可编辑代码及其编码能力；外层探索固定 & 基准分数 & $k$；元层另辨 & 自修改实例；不代表整个修改机制均可变 \\
JitRL\cite{jitrl-2601} & 运行时 logit 偏置（非持久组件），$M$ 为非参库 & LLM轨迹评分 & $e$ + $k$（$M$） & 不更新基座参数的对照；仍有经验与推理费 \\
AgeMem\cite{agemem-2601} & 训练更新$\theta$，任务内操作$M$与工作上下文 & 任务、上下文与记忆复合奖励 & 训练与任务内分开 & 记忆操作学习化的实例；第 6 章 \\
MEM1\cite{mem1-2506} & $\theta$（PPO）；任务内内部状态 & 环境结果（可验证奖励） & $k$（$\theta$） & 学习状态管理；不是部署期持续训练 \\
Mem-$\alpha$\cite{memalpha-2509} & $\theta$中的写入策略；结构固定 & QA、格式、压缩与合法性加权奖励 & $k$（写入策略） & 记忆构建策略学习；第6章 \\
MemGen\cite{memgen-2509} & 触发器与编织器；任务内潜记忆 & 依任务而定 & 训练与任务内分开 & 不直接读出语义；仍可做行为和消融检验 \\
ACE\cite{ace-2510} & $P$（上下文条目） & 执行及可得的任务反馈 & $k$（$P$） & 第7章；需区分有标签与无标签设置 \\
GEPA\cite{gepa-2507} & $P$（提示） & 评价指标与轨迹反馈 & $k$（$P$） & 第 7 章；rollout 数之比不是成本之比 \\
\end{longtable}
\endgroup

\begin{sloppypar}
看这张表时要抓住的只有一条：\textbf{“改什么”这一栏决定了它需要什么样的检验}。
改 $\theta$ 的方法必须回答训练数据与检验数据的关系；改 $M$ 的方法必须回答条目何时被读、何时被淘汰；
改 $G$ 的方法必须回答搜索用的验证集与最终报告集是否分开；改修改器的方法还必须单独检验元层本身。
把这一栏空着，后面所有数字都只能读作“某个整体变好了”，读不出是哪一部分在起作用。
\end{sloppypar}

% ---------------- E.5 ----------------
\section{易错处}

\begin{pitfallbox}
\textbf{反例：把层数不同的两个“记忆”并列比较。}
ReAct 的任务内推理与动作，和 Reflexion 的跨 trial 反思保留，具有不同的信息路径。后者是否跨新任务仍取决于具体复位协议。
若把两者的成功率直接相减，差异可能同时包含推理方式、反馈、重试次数与记忆保留规则，不能仅归因于持久化。
本书的做法是给每个方法先填“落在哪一层”（表 \ref{tab:appE-methods} 第四列），
层级不同的方法也可以在共同任务和预算下比较；需要说明该差代表整体方案，不能未经消融便归因于持久性。
\end{pitfallbox}

\begin{pitfallbox}
\textbf{反例：把“验证器”当成“验收器”的同义词。}
本书的 $\mathcal V$ 含任务划分、验收器、最终测试集三部分。
只冻结验收器而让系统参与划分测试集，等于把最终测试集交给了被评对象——
自生成题的成绩仍有描述意义，但不能直接代表独立目标分布上的收益。第 3 章的封存集纪律正是为这一处设的。
\end{pitfallbox}

\begin{pitfallbox}
\textbf{反例：用“演化成本”只算训练算力。}
$C_{\text{evolve}}$ 包含构建、验证、回归三部分的全部资源。
完整方案还需另列部署与推理费用。E03 拆分 B 把同一份拟合求解器的记忆交给原读取者和更换后的读取者。两臂都在同一组 \measured{16} 道题上各运行 \measured{5} 次，回合计时和分别为 \measured{374.2} 秒与 \measured{2084.9} 秒（§\ref{sec:appF-cross}）。
整次 E03 续跑跨日，这两臂的留存轨迹都始于 9 月 26 日（UTC），先用原读取者，再换模型。模型与运行时段一起变化，计时也只覆盖所记录的回合。这组数值描述两种读取条件下的耗时，无法单独识别更换模型的作用或给出完整演化费用。
\end{pitfallbox}

\openquestion{对照表只提供定位入口，不是效果排行榜。第10—12章的参数学习、课程与自修改方法，应结合各章的具体实现边界与来源索引阅读；缺失训练成本、对照或原始日志的字段仍需保留。}
